% Agent4IR '26 Workshop Paper — KDD 2026
% Reviewed and strengthened version
\documentclass[sigconf,nonacm]{acmart}

\usepackage{booktabs}
\usepackage{multirow}

% Remove ACM copyright for workshop submission
\setcopyright{none}
\settopmatter{printacmref=false}
\renewcommand\footnotetextcopyrightpermission[1]{}
\pagestyle{plain}

\begin{document}

\title{Biomedical QA Agents as Evidence Mediators: Agentic Retrieval, DSPy Baselines, and the Moral Cost of Binary Answers}

\author{Anonymous Author(s)}
\affiliation{}

\begin{abstract}
Biomedical question answering is often treated as a technical pipeline: retrieve PubMed abstracts, prompt a language model, and score the answer. We argue that this framing misses the central problem. Biomedical QA is evidence labor performed under time pressure, where query choices, missing items, and premature yes/no closure shape what clinicians and researchers may notice. We present a bounded agentic RAG system for BioASQ-style QA that lets an LLM decompose questions, search hybrid dense/sparse indexes, judge evidence sufficiency, reformulate queries, rerank passages, generate type-specific answers, verify claims, and merge consensus outputs. We compare this domain-governed agent with DSPy baselines---zero-shot, RAG, optimized RAG, MultiHop RAG, ReAct RAG, and optimized ReAct RAG---built over the same PubMed retrieval substrate. On BioASQ-style validation, the strongest hybrid agent reaches 76.0 overall score, compared with 71.5 for a simpler agentic FTS5 system and 45.5 for the best DSPy baseline. Gains are concentrated in factoid and list questions, where static or generic agentic programs frequently miss distributed evidence. Our novelty is not a new retriever or language model, but a sociotechnical account of agentic IR: agents are useful when they make iterative evidence work inspectable, but dangerous when they make binary biomedical answers feel morally complete.
\end{abstract}

\maketitle

\section{Introduction}

Biomedical literature search is a high-stakes form of information work. PubMed reports more than 40 million biomedical citations, while BioASQ asks systems to retrieve articles and snippets and to produce exact and ideal answers for expert-written biomedical questions~\cite{bioasq_task14b, bioasq_challenge}. The challenge is not simply scale. Biomedical search is shaped by synonymy, abbreviations, entity exactness, changing evidence, and differences between what a benchmark can score and what a clinician or researcher can responsibly use. Prior work on clinical information seeking shows that medical professionals encounter frequent questions and face barriers of time, access, and search skill~\cite{bastregari2020}. AI-based literature tools may help, but only if they do more than produce fluent summaries~\cite{pubmed_beyond}.

Standard retrieval-augmented generation (RAG) performs retrieval once and then generates an answer~\cite{lewis2020rag}. This is attractive because it is simple and cheap, but it is a poor model of expert evidence seeking. A biomedical expert often tries several query formulations, checks whether evidence is sufficient, separates exact entity extraction from narrative synthesis, and revises conclusions when contradictions appear. We therefore study an agentic RAG system in which an LLM controls a constrained search-and-reasoning loop. The agent is not allowed to ``practice medicine''; it is allowed to perform bounded evidence mediation: ask better retrieval questions, inspect evidence, report uncertainty, and produce answer forms required by BioASQ.

This paper makes three claims. First, in BioASQ-style QA, run-time control over the evidence trajectory matters more than prompt optimization alone for list and factoid questions. Second, not all ``agentic'' systems are equivalent: a generic ReAct loop differs materially from a domain-governed agent that has explicit sufficiency checks, answer-type policies, verification, and list-recall passes. Third, the strongest contribution of biomedical QA agents is sociotechnical rather than purely algorithmic. Agents are worth building when they expose and improve evidence labor; they are morally hazardous when they compress uncertain biomedical evidence into unexamined yes/no outputs.

\section{Task, Data, and Evaluation Design}

We use BioASQ Task 13b as the target task. The task contains English biomedical questions with gold articles, snippets, exact answers, and ideal answers constructed by biomedical experts. Task 13b has four answer types: yes/no, factoid, list, and summary. Phase~A asks systems to return up to 10 relevant PubMed articles and snippets; Phase~A+ asks for exact and ideal answers from system-retrieved evidence before gold evidence is released; Phase~B supplies gold documents and snippets and asks systems for exact and ideal answers~\cite{bioasq_task14b}. The organizers focus on PubMed titles and abstracts, not PubMed Central full text, and the 2026 Phase~A designated source is the PubMed Annual Baseline Repository~\cite{pubmed_home}.

This task design matters for both technical and social reasons. Factoid and list questions require exact entities rather than paraphrases; BioASQ explicitly discourages submitting synonyms for exact factoid and list answers~\cite{bioasq_task14b}. Summary questions require paragraph-sized ideal answers, and yes/no questions require binary exact answers even though the evidence behind the answer may be conditional. Thus BioASQ is not one task but four epistemic regimes: entity ranking, entity set completion, binary closure, and narrative synthesis. A system can do well on one regime and fail another. This is why we report type-level results instead of treating overall accuracy as sufficient.

Our development corpus includes over 5,200 historical BioASQ training questions with gold evidence and answers; the official Task 13b development file contains 5,389 questions~\cite{bioasq_task14b}. For retrieval experiments, the 2026 PubMed abstract-only baseline contains 28,336,648 documents in our index. We report internal validation results on a held-out 85-question scored slice for zero-shot (26 factoid, 17 yes/no, 23 list, 19 summary) and broader internal run summaries for the agentic and DSPy systems. Metrics follow the BioASQ spirit: yes/no macro-F1, factoid MRR or exact-match-style accuracy, list F1, and ideal-answer semantic evaluation; Phase~A retrieval uses MAP and F-measure conventions~\cite{bioasq_eval}. Because some recorded runs used count-based correctness and others used the unweighted mean of task metrics, Table~\ref{tab:results} marks this distinction explicitly.

\section{Systems Compared}

\subsection{Agentic Evidence-Seeking Pipeline}

Our primary system is a bounded ReAct-style biomedical evidence agent~\cite{yao2023react}. The central design decision is to treat biomedical question answering not as a simple retrieve-then-generate operation, but as a controlled evidence-seeking process in which the model must repeatedly decide what it knows, what it does not yet know, and what evidence would be sufficient to justify an answer. This distinction matters in BioASQ-style QA because many questions are not simply lexical lookup problems. A question may require mapping lay and technical terminology, linking a gene or pathway to a disease phenotype, distinguishing mechanistic evidence from clinical evidence, or aggregating many entities into a list answer. In such cases, a fixed retrieval pass can fail silently: it may retrieve plausible documents, generate a fluent answer, and provide no opportunity to notice that the evidence was incomplete. Our agentic pipeline instead makes evidence acquisition a runtime policy. The system can search, inspect, reformulate, and only then answer.

The agent is implemented as an eight-stage loop: \emph{think, search, evaluate, rerank, answer, verify, list review}, and \emph{consensus}. In the \emph{think} stage, a local instruction-tuned 31B model first identifies the BioASQ question type---factoid, yes/no, list, or summary---and predicts whether the question is likely to require multi-hop reasoning. When multi-hop reasoning is detected, the model decomposes the original question into two to four evidence-seeking subquestions. For example, a question about the cell state associated with \emph{NF1} loss in glioblastoma requires evidence about the role of NF1, the relevant glioblastoma cell state taxonomy, and the empirical association between loss of NF1 and that state. Rather than relying on one broad query, the agent creates several targeted searches that separate these evidentiary needs. This decomposition is not intended to make the agent autonomous in an open-ended sense; it is a practical mechanism for exposing the intermediate assumptions that a conventional RAG system would hide inside a single prompt.

In the \emph{search} stage, each query is submitted to a hybrid biomedical retrieval stack. Dense retrieval uses NeuML's PubMedBERT embedding model, which maps biomedical text into a 768-dimensional vector space~\cite{neuml_embeddings}. To make retrieval over the PubMed 2026 abstract-only baseline computationally feasible, the dense index is built with FAISS IVF-Flat using 8,192 Voronoi cells ($\texttt{nlist}=8192$) and probing 91 cells at search time ($\texttt{nprobe}=91$)~\cite{faiss}. This approximate nearest-neighbor structure gives the agent fast access to semantic matches, including cases where the wording of the question differs from the wording of the article. Sparse retrieval uses BM25 with default parameters~\cite{robertson2009bm25}. We keep BM25 because biomedical questions often contain high-value exact tokens---gene symbols, drug names, diseases, variants, and molecular targets---that should not be softened away by embedding similarity. The dense and sparse rankings are merged using reciprocal rank fusion (RRF), with semantic retrieval weighted at 0.6 and sparse retrieval weighted at 0.4~\cite{cormack2009rrf}. This design gives the agent both conceptual recall and lexical precision.

The \emph{evaluate} stage is the first point where the system departs substantially from ordinary RAG. After the initial hybrid retrieval, the model reads the top evidence and explicitly judges whether the available passages are sufficient to answer the question. If the answer is \textsc{insufficient}, the agent generates three new search queries and returns to retrieval. These new queries are not minor paraphrases of the first query; they are prompted to target missing mechanisms, missing entities, contradictory evidence, or under-specified biomedical terms. This loop is bounded to a maximum of four retrieval iterations. This bound is important both computationally and epistemically. Computationally, it prevents unbounded tool use. Epistemically, it frames agency as disciplined uncertainty management rather than unconstrained autonomy. The agent is useful precisely because it can say, in effect, ``the current evidence is not enough,'' and then perform a targeted search to reduce that uncertainty.

After iterative search, the \emph{rerank} stage applies a biomedical cross-encoder, NeuML/biomedbert-base-reranker, to the accumulated candidate passages~\cite{neuml_reranker}. Unlike the FAISS retriever, which compares precomputed vector representations, the cross-encoder processes the query and each passage jointly with full cross-attention. This allows the system to rescore evidence according to the specific question being asked rather than only according to global embedding proximity. The reranked evidence set is then compressed into the passages used for answer generation. This division of labor is deliberate: FAISS and BM25 provide broad recall at scale, while the cross-encoder supplies a more expensive but more precise judgment over a smaller candidate pool.

The \emph{answer} stage uses type-specific prompting rather than a single universal generation prompt. For factoid questions, the model is instructed to copy the exact biomedical entity or term from the evidence and avoid paraphrase. This reduces the common failure mode in which a language model produces a semantically related but evaluation-incompatible answer. For yes/no questions, the model is prompted adversarially: it must look for evidence supporting both affirmative and negative conclusions before deciding. This is technically useful for avoiding one-sided retrieval bias, but it is also central to our broader argument. A yes/no biomedical answer is a form of moral compression: a complex and uncertain evidence base is collapsed into a binary output that may later be read as a recommendation. Requiring the model to inspect both sides makes that compression more accountable. For list questions, the model proceeds passage by passage and extracts all candidate items before normalizing and deduplicating. For summary questions, the model writes a concise ideal answer grounded in the retrieved snippets rather than a free-form biomedical explanation.

The final three stages are designed to reduce unsupported fluency. In \emph{verify}, the model checks the draft answer against the evidence and identifies claims that are missing, contradicted, or not directly supported. When unsupported claims are found, the answer is revised or shortened rather than expanded. In \emph{list review}, used only for list questions, the model performs a second extraction pass asking what entities may have been missed in the first pass. This stage is motivated by the observation that list answers are recall-sensitive: a system can identify several correct items and still fail because it stops too early. The second pass typically recovers additional candidates that are then normalized and merged. Finally, in \emph{consensus}, the system combines three independently generated answers produced with low but distinct temperatures. For factoid and yes/no questions, we use majority voting; for list questions, we use a union-and-normalization procedure; for ideal answers, we select a stable median-length synthesis that preserves evidence coverage without rewarding verbosity.

We therefore use the term \emph{agentic} in a deliberately bounded sense. The system is not an unconstrained clinical actor, and it is not allowed to invent tools, browse arbitrary sources, or make decisions outside the evidence context. Its agency is located in the retrieval-and-reasoning loop: deciding when evidence is insufficient, selecting new searches, comparing competing passages, and revising answers when support is weak. This bounded form of agency is the methodological core of the paper. It converts biomedical QA from a one-shot prediction task into an auditable sequence of evidence labor. From a computational social science perspective, this matters because the social risk of biomedical AI is not only that a model may be wrong; it is that a wrong answer may appear frictionless, authoritative, and detached from the labor required to justify it. Our pipeline attempts to reintroduce that labor into the system design by making uncertainty, search, verification, and answer compression explicit parts of the method.

\subsection{DSPy Baselines and Optimization}

We compare against a DSPy suite because it represents a strong alternative view: rather than handcrafting prompts, express language-model programs as modules and optimize their instructions and demonstrations~\cite{dspy}. We evaluate zero-shot, RAG, optimized RAG, MultiHop RAG, ReAct RAG, and optimized ReAct RAG. All non-zero-shot baselines use \texttt{gpt-4o-s-120b} as generation model and, for optimization, as teacher; the zero-shot baseline uses Gemma 4 31B-IT. OpenAI describes \texttt{gpt-4o-s-120b} as an open-weight model designed for reasoning and tool-use deployment~\cite{openai_gpt4os}.

The DSPy baselines use the same retrieval substrate as the agentic system: NeuML PubMedBERT dense retrieval, BM25 sparse retrieval, RRF fusion with default hyperparameters, FAISS IVF-Flat over 28,336,648 PubMed abstract-bearing documents, and the NeuML biomedical cross-encoder reranker. The basic RAG program retrieves 1,000 documents, reranks, and passes the top 30 documents to a DSPy \texttt{ChainOfThought} module with the signature \texttt{context, question -> response}. MultiHop RAG performs four hops; each hop retrieves 1,000 documents, reranks, summarizes the top 10 into notes, and uses accumulated notes to refine the next query. ReAct RAG uses DSPy's built-in \texttt{dspy.ReAct} module with the signature \texttt{question -> response}; the retrieval pipeline is exposed as a search tool, each tool call returns the top 10 reranked documents, and the agent may take up to 20 iterations.

For all optimized DSPy settings, we use MIPROv2, which jointly proposes and optimizes instructions and few-shot examples using program- and data-aware search~\cite{miprov2, khattab2024optimizing}. Optimization draws up to 200 training examples, 300 development examples, and 500 test examples at random from \texttt{rag-mini-bioasq}, a Hugging Face dataset derived from BioASQ Task 11b training data~\cite{rag_datasets}. Only question and answer fields are used during optimization, not gold snippets. We also use DSPy's out-of-the-box SemanticF1 metric in decomposition mode, with Kimi K2.6 as the evaluator model~\cite{dspy_semanticf1, dspy_eval}. This setup makes the comparison informative but not perfectly symmetric: DSPy optimizes prompting/program behavior, while our agent adds biomedical control policies and answer-type obligations.

\section{Findings}

\begin{table}[t]
\caption{Internal BioASQ-style validation results. Values are percentages except where noted. ``Summary'' is available only for the zero-shot count-based run. The zero-shot overall is question-level correctness over 85 examples; all other overall values are the unweighted mean of Factoid, Yes/No, List, and SemanticF1.}
\label{tab:results}
\small
\begin{tabular}{lccccc}
\toprule
\textbf{System} & \textbf{Fac.} & \textbf{Y/N} & \textbf{List} & \textbf{SemF1} & \textbf{Ovr.} \\
\midrule
Zero-shot Gemma 4 31B & 8.0 & 94.0 & 0.0 & 41.6 & 35.7$^\dagger$ \\
\midrule
Agentic SQLite FTS5 & 62.0 & 88.0 & 87.0 & 49.0 & 71.5 \\
Agentic Hybrid FAISS+BM25+IVF & 69.0 & 88.0 & \textbf{96.0} & 51.0 & \textbf{76.0} \\
Agentic SQLite FTS5, COT prompt & 50.0 & 82.0 & 70.0 & 49.0 & 62.75 \\
\midrule
DSPy RAG & 12.0 & 82.0 & 4.0 & 49.0 & 36.75 \\
DSPy RAG + MIPROv2 & 27.0 & 94.0 & 13.0 & 48.0 & 45.5 \\
DSPy MultiHop RAG & 12.0 & 76.0 & 0.0 & 49.0 & 34.25 \\
DSPy MultiHop RAG + MIPROv2 & 27.0 & 88.0 & 13.0 & 48.0 & 44.0 \\
DSPy ReAct RAG & 8.0 & 76.0 & 0.0 & 55.0 & 34.75 \\
DSPy ReAct RAG + MIPROv2 & 31.0 & \textbf{94.0} & 0.0 & \textbf{55.0} & 45.0 \\
\bottomrule
\end{tabular}
\vspace{-2mm}
\end{table}

Table~\ref{tab:results} shows three patterns. First, the best performance comes from a domain-governed agent, not from a generic ReAct loop or from optimized static RAG. The hybrid agent reaches 76.0 overall, improving over the SQLite FTS5 agent by 4.5 points and over the best DSPy baseline by 30.5 points. The improvement is largest for list questions: the hybrid agent reaches 96.0, while DSPy RAG variants range from 0.0 to 13.0. List questions are recall-sensitive and distributed across evidence; this is exactly where query reformulation, evidence sufficiency checks, and a second missed-item pass should help.

Second, prompt/program optimization improves some baselines but does not solve evidence completeness. MIPROv2 raises DSPy RAG factoid accuracy from 12.0 to 27.0 and yes/no from 82.0 to 94.0, yet list performance remains low. Optimized ReAct reaches 94.0 on yes/no and 55.0 SemanticF1, but 0.0 on list. This suggests that optimization can make answer style and binary decisions more benchmark-aligned without guaranteeing that the system has found the full evidence set.

Third, yes/no and summary-style results can be misleadingly comforting. The zero-shot Gemma baseline achieves 16/17 yes/no and 19/19 summary correctness under the recorded count-based evaluation, but only 2/26 factoid and 0/23 list. This is not a trivial failure mode. It means a model can sound useful when the output is binary or narrative while failing at the entity-level evidence work that biomedical users often need. From a technical perspective, this motivates type-specific routing and verification. From a sociotechnical perspective, it warns against evaluating biomedical agents mainly by fluent answers or aggregate scores.

\section{Discussion}

\subsection{Why Agents at All?}

The question ``Why do we need agents for biomedical QA?'' should not be answered with model excitement. We need agents only if they change the organization of evidence work in a desirable way. Biomedical search is a form of invisible labor: users translate vague questions into database queries, remember synonyms, infer which abstracts are worth opening, decide when evidence is sufficient, and summarize findings for someone else~\cite{suchman1995}. Static RAG automates only one slice of this labor: an initial query and a top-$k$ cutoff. An agentic retriever can instead make this labor explicit: the generated subqueries, failed searches, sufficiency judgments, reranker decisions, and verification edits become inspectable artifacts.

This is the sociotechnical novelty of our approach. The individual components---BM25, FAISS, PubMedBERT, RRF, cross-encoder reranking, ReAct, and prompt optimization---are not new. The contribution is to treat the \emph{evidence trajectory} as the unit of design and evaluation. In our results, systems that share a retrieval substrate behave very differently depending on the control policy around that substrate. Generic DSPy ReAct has tool use, but lacks biomedical obligations such as adversarial yes/no checking and list recall review. The hybrid agent's gains therefore argue for governed agency: not more autonomy, but more explicit procedural commitments about how biomedical evidence must be searched, checked, and represented.

\subsection{The Moral Cost of Yes/No}

Yes/no biomedical questions create a moral compression problem. A binary answer is useful for evaluation and sometimes useful in practice, but it collapses evidence quality, population heterogeneity, clinical context, and uncertainty into a single token. Our results show why this matters. The systems with the best yes/no scores are not necessarily the systems with the best evidence completeness. A tool that answers ``yes'' correctly but cannot retrieve the relevant entities, mechanisms, or exceptions may still encourage automation bias, a known concern for AI-based clinical decision support~\cite{khera2023, sutton2020}.

The agent's adversarial yes/no procedure is a technical response to this moral problem. Before answering, the prompt requires evidence for both ``yes'' and ``no'' and asks the model to state why one side is better supported. This does not make the system morally safe. It does, however, change binary closure from a generation habit into an accountable step. The right interface should preserve this distinction: the user should see not only ``yes'' or ``no,'' but also the evidence trail, the rejected alternative, and any insufficiency flag. In biomedical IR, a correct binary label without evidence provenance is often not enough to be trustworthy.

\subsection{Temporal Evidence and Biomedical Drift}

A further reason biomedical QA cannot be treated as ordinary document retrieval is that biomedical knowledge is temporally unstable. PubMed is cumulative: older abstracts remain searchable even when later trials, reviews, guidelines, or failures to replicate have changed the interpretation of the original claim. In a benchmark, a gold snippet can be treated as stable evidence for a question. In clinical and research practice, however, evidence has a life cycle. A paper may be preliminary, superseded, contradicted, narrowed to a particular population, or reinterpreted after new methods appear. A retrieval system that ignores this temporal structure can answer from evidence that is relevant in the lexical sense but stale in the epistemic sense.

This creates a distinct failure mode for biomedical RAG. The retriever is usually optimized to find semantically similar passages, not to determine whether those passages still represent the best available knowledge. Dense retrieval can surface an older mechanistic study because it is conceptually close to the question; BM25 can surface an older abstract because it contains the exact gene, drug, or disease name. Neither signal by itself tells the model whether the claim has been updated, contradicted, or absorbed into a newer consensus. The language model may then perform what looks like synthesis while actually flattening time: a 2008 association study, a 2017 animal model, a 2022 randomized trial, and a 2025 review can be presented as if they occupy the same evidentiary plane.

This is especially dangerous for yes/no questions. A binary answer can make an old claim feel current. If an agent retrieves older affirmative evidence but misses newer negative or more qualified evidence, the output may be technically supported by a snippet while socially misleading. The moral problem is not only hallucination; it is temporal misrepresentation. A model may cite a real paper and still give a bad answer because it fails to ask whether that paper should still govern the conclusion. In biomedical settings, correctness is therefore partly temporal: \emph{when} the evidence was produced, \emph{what kind} of evidence it is, and \emph{whether} later evidence has changed the claim are all part of the answer.

Our current system does not fully solve temporal evidence drift. The BioASQ-style setting rewards retrieval of relevant PubMed abstracts and exact answer production, not longitudinal evidence surveillance. Nevertheless, the agentic architecture is better suited to this problem than a one-shot retriever because it already contains explicit stages for sufficiency judgment, query reformulation, verification, and consensus. These stages can be extended into temporal safeguards: date-aware reranking, contradiction-seeking queries, preference for recent high-level evidence when appropriate, and abstention when retrieved evidence is too old or too mixed to justify binary closure. From a computational social science perspective, this matters because biomedical authority is not only distributed across documents; it is distributed across time. Responsible agentic IR should therefore make the temporality of evidence visible rather than allowing old and new claims to collapse into the same fluent answer.

\subsection{Accountability Without the Moral Crumple Zone}

Agentic systems also risk creating what Elish calls a moral crumple zone: humans absorb blame for failures of automated systems that they could not realistically understand or control~\cite{elish2019}. A clinician shown a polished answer may be formally ``in the loop'' while having little visibility into query reformulations, omitted abstracts, or ranking errors. Our design principle is therefore not human replacement but accountable handoff. The agent should provide provenance for retrieved documents and snippets, record its search path, expose sufficiency decisions, and abstain or flag uncertainty when evidence is weak. These are not only user-interface features; they are IR requirements for responsible biomedical agents.

This discussion also reframes evaluation. A four-column score table is necessary but insufficient. We should evaluate whether traces help humans detect errors, whether list omissions cluster around rare diseases or under-studied populations, whether generated queries systematically privilege well-indexed terminology, and whether users over-trust high-confidence binary answers. In this sense, agent logs are computational social science data: they reveal how an automated system operationalizes relevance, uncertainty, and biomedical authority. WHO guidance argues that AI for health must center ethics, human rights, transparency, responsibility, and inclusiveness~\cite{who2021}. For agentic biomedical QA, those values become concrete design obligations: bounded tool use, evidence provenance, calibrated abstention, answer-type-specific safeguards, and human review that is meaningful rather than ceremonial.

\section{Limitations and Conclusion}

Our results are internal validation results, not official BioASQ leaderboard scores. The zero-shot row and the other run summaries use different aggregate conventions, so conclusions should focus on type-level patterns rather than a single overall number. The corpus is title/abstract-only, excluding full-text evidence. DSPy optimization used question-answer pairs from \texttt{rag-mini-bioasq}, not gold snippets, which may disadvantage evidence-grounded behavior. Finally, logged agent traces are not automatically usable by clinicians or researchers; future work must test whether traces actually improve human judgment.

We conclude that biomedical QA agents should be framed as evidence mediators, not medical authorities. The technical finding is that a domain-governed agentic loop improves recall-sensitive and exact-answer regimes over static and generic agentic baselines. The computational social science finding is that the value of agents lies in making evidence labor visible and governable. The central question is not only whether the answer is correct, but whether the path to the answer can be inspected, challenged, and responsibly handed back to human experts.

\bibliographystyle{ACM-Reference-Format}
\begin{thebibliography}{26}

\bibitem{bastregari2020}
S.~Bastregari et~al. 2020. Clinical information seeking behavior of physicians: A systematic review. \emph{International Journal of Medical Informatics} 139 (2020), 104145.

\bibitem{bioasq_task14b}
BioASQ. 2026. BioASQ Task 14b Guidelines. \url{http://participants-area.bioasq.org/general_information/Task14b/}.

\bibitem{bioasq_challenge}
BioASQ. 2026. The BioASQ Challenge. \url{https://www.bioasq.org/}.

\bibitem{cormack2009rrf}
G.~V. Cormack, C.~L.~A. Clarke, and S.~Buettcher. 2009. Reciprocal Rank Fusion Outperforms Condorcet and Individual Rank Learning Methods. In \emph{Proc.\ SIGIR}. 758--759.

\bibitem{faiss}
M.~Douze et~al. 2024. The Faiss Library. \emph{arXiv:2401.08281}.

\bibitem{miprov2}
DSPy. 2026. MIPROv2 Optimizer. \url{https://dspy.ai/api/optimizers/MIPROv2}.

\bibitem{dspy_semanticf1}
DSPy. 2026. SemanticF1 Evaluation. \url{https://dspy.ai/api/evaluation/SemanticF1}.

\bibitem{elish2019}
M.~C. Elish. 2019. Moral Crumple Zones: Cautionary Tales in Human-Robot Interaction. \emph{Engaging Science, Technology and Society} 5 (2019), 40--60.

\bibitem{pubmed_beyond}
Yu~Gu et~al. 2021. Domain-Specific Language Model Pretraining for Biomedical Natural Language Processing. \emph{ACM Trans.\ Computing for Healthcare} 3, 1 (2021), 1--23.

\bibitem{pubmed_home}
Qiao~Jin, R.~Leaman, and Z.~Lu. 2024. PubMed and beyond: biomedical literature search in the age of artificial intelligence. \emph{eBioMedicine} 100 (2024), 104988.

\bibitem{dspy}
O.~Khattab et~al. 2023. DSPy: Compiling Declarative Language Model Calls into Self-Improving Pipelines. \emph{arXiv:2310.03714}.

\bibitem{khera2023}
R.~Khera, M.~Simon, and J.~S. Ross. 2023. Automation Bias and Assistive AI: Risk of Harm from AI-Driven Clinical Decision Support. \emph{JAMA} 330, 23 (2023), 2255--2257.

\bibitem{lewis2020rag}
P.~Lewis et~al. 2020. Retrieval-Augmented Generation for Knowledge-Intensive NLP Tasks. In \emph{NeurIPS}. 9459--9474.

\bibitem{openai_gpt4os}
Moonshot AI. 2026. Kimi K2.6: Advancing Open-Source Coding. \url{https://www.kimi.com/blog/kimi-k2-6}.

\bibitem{neuml_embeddings}
NeuML. 2026. NeuML/pubmedbert-base-embeddings. \url{https://huggingface.co/NeuML/pubmedbert-base-embeddings}.

\bibitem{neuml_reranker}
NeuML. 2026. NeuML/biomedbert-base-reranker. \url{https://huggingface.co/NeuML/biomedbert-base-reranker}.

\bibitem{khattab2024optimizing}
K.~O.~Ong et~al. 2024. Optimizing Instructions and Demonstrations for Multi-Stage Language Model Programs. \emph{arXiv:2406.11695}.

\bibitem{rag_datasets}
RAG Datasets. 2024. rag-datasets/rag-mini-bioasq. \url{https://huggingface.co/datasets/rag-mini-bioasq}.

\bibitem{robertson2009bm25}
S.~Robertson and H.~Zaragoza. 2009. The Probabilistic Relevance Framework: BM25 and Beyond. \emph{Foundations and Trends in Information Retrieval} 3, 4 (2009), 333--389.

\bibitem{suchman1995}
S.~L. Star and A.~Strauss. 1999. Layers of Silence, Arenas of Voice: The Ecology of Visible and Invisible Work. \emph{Computer Supported Cooperative Work} 8, 1--2 (1999), 9--30.

\bibitem{sutton2020}
R.~T. Sutton et~al. 2020. An overview of clinical decision support systems: benefits, risks, and strategies for success. \emph{NPJ Digital Medicine} 3 (2020).

\bibitem{mark2025}
M.~Sujan et~al. 2025. Exploring the Risks of Automation Bias in Healthcare AI. \emph{Intelligence-Based Medicine} 10 (2025), 100194.

\bibitem{who2021}
World Health Organization. 2021. Ethics and Governance of Artificial Intelligence for Health: WHO Guidance. \url{https://www.who.int/publications/i/item/9789240029200}.

\bibitem{bioasq_eval}
G.~Balikas et~al. 2013. Evaluation Framework Specifications. Technical Report D4.1, BioASQ Deliverable.

\bibitem{yao2023react}
S.~Yao et~al. 2023. ReAct: Synergizing Reasoning and Acting in Language Models. In \emph{ICLR}.

\bibitem{yu2007}
H.~Yu et~al. 2007. Development, implementation, and a cognitive evaluation of a definitional question answering system for physicians. \emph{Journal of Biomedical Informatics} 40, 3 (2007), 236--251.

\end{thebibliography}

\end{document}
